\chapter{Couche 3 — Les machines louées}

C'est le cœur du produit. Sept briques pour fabriquer ce que l'utilisateur
achète : une machine Linux à lui, joignable de trois façons.

% -----------------------------------------------------------------------------
\outil{Docker}{brique 15 · la conteneurisation}

\quoi
Un moteur qui exécute des processus isolés les uns des autres, en réutilisant le
noyau de la machine hôte. Un \emph{conteneur} se croit seul au monde, mais ne
porte pas de système d'exploitation complet.

\pourquoi
La vraie alternative était la machine virtuelle --- VirtualBox, KVM. Une VM
offre une isolation plus forte, puisqu'elle a son propre noyau. Deux chiffres
ont tranché :

\begin{center}
\small
\begin{tabularx}{0.9\textwidth}{@{}lXX@{}}
\toprule
& \textbf{Machine virtuelle} & \textbf{Conteneur} \\
\midrule
Démarrage & 30 à 60 secondes & moins d'une seconde \\
Mémoire minimale & plusieurs centaines de Mo & celle des processus \\
Densité sur un serveur & quelques dizaines & plusieurs centaines \\
\bottomrule
\end{tabularx}
\end{center}

Pour une location \textbf{à la minute}, faire attendre l'utilisateur une minute
est absurde. Le sujet tranche d'ailleurs explicitement : « des conteneurs dans
notre cas ».

\begin{notion}{Ce qu'un conteneur est réellement}
Trois mécanismes du noyau Linux, pas un objet magique :

\textbf{Les espaces de noms} cloisonnent ce que le processus \emph{voit} : son
propre arbre de processus, ses propres interfaces réseau, sa propre vision du
disque. Le port 22 du conteneur n'est pas le port 22 de l'hôte.

\textbf{Les groupes de contrôle} limitent ce qu'il \emph{consomme} : mémoire,
processeur, nombre de processus.

\textbf{Les capacités} découpent les pouvoirs de root en une quarantaine de
droits indépendants, qu'on accorde un par un.
\end{notion}

\role
Six images, et une commande imposée par le sujet pour créer chaque machine :

\begin{lstlisting}[language=bash]
docker run -d --restart=always -p <port>:22 <image>
\end{lstlisting}

Le projet y ajoute le durcissement : \cmd{--cap-drop ALL} puis les dix capacités
strictement nécessaires, \cmd{no-new-privileges}, une limite de processus, de
CPU, et la mémoire en réservation \emph{et} en limite.

\apport
La plateforme sait livrer une machine en quinze secondes. C'est la brique qui
rend le produit possible.

% -----------------------------------------------------------------------------
\outil{Dockerfile}{brique 16 · la recette des images}

\quoi
Un fichier texte qui décrit, instruction par instruction, comment construire
une image. Chaque instruction crée une \emph{couche} empilée sur la précédente.

\pourquoi
On pourrait créer une image en lançant un conteneur, en l'installant à la main,
puis en figeant le résultat. Personne ne saurait plus ce qu'il contient, ni
comment le reproduire.

Le Dockerfile est l'infrastructure comme code appliquée aux images : versionné,
relu, et vérifiable par un analyseur.

\role
Quatre fichiers pour six images. Le mode bureau n'est pas un fichier séparé :
c'est un paramètre de construction.

\begin{lstlisting}[language=dockerfile]
FROM ubuntu:22.04
ARG DESKTOP=0
RUN apt-get update \
    && apt-get install -y --no-install-recommends openssh-server supervisor \
    && rm -rf /var/lib/apt/lists/*
COPY entrypoint.sh /usr/local/bin/entrypoint.sh
CMD ["/usr/local/bin/entrypoint.sh"]
\end{lstlisting}

\begin{piege}{Nettoyer dans une autre couche ne sert à rien}
Si l'on installe des paquets dans un \cmd{RUN} et qu'on efface le cache d'apt
dans le \cmd{RUN} suivant, le cache reste dans l'image : une couche ne peut pas
alléger celle du dessous. Il faut enchaîner les deux dans la \emph{même}
instruction. Sur nos images, cela représente une cinquantaine de mégaoctets.
\end{piege}

\begin{piege}{Un mot de passe par défaut dans l'image}
Notre première version contenait \cmd{ENV ROOT\_PASSWORD=insacloud}. L'analyseur
\cmd{hadolint} l'a signalé, et il avait raison : ce mot de passe aurait été le
même sur \textbf{toutes} les machines louées, et visible par quiconque fait
\cmd{docker history}. L'entrypoint en tire désormais un au hasard si la
plateforme n'en fournit pas.
\end{piege}

\apport
Six images reproductibles, construites par la chaîne d'intégration à chaque
envoi, et analysées avant d'être publiées.

% -----------------------------------------------------------------------------
\outil{supervisord}{brique 17 · plusieurs services dans un conteneur}

\quoi
Un gestionnaire de processus : il en démarre plusieurs, les surveille, et
relance celui qui meurt.

\pourquoi
C'est la brique qui résout une contradiction. Un conteneur exécute \emph{un}
processus, et s'arrête quand celui-ci se termine. Or une machine louée doit faire
tourner cinq services : SSH, le terminal web, le serveur d'affichage, le bureau,
et la passerelle VNC.

Les alternatives : lancer un script shell qui démarre tout en arrière-plan ---
mais alors rien ne surveille, et un service mort le reste. Ou \cmd{systemd} dans
le conteneur --- lourd, et il réclame des privilèges qu'on vient justement de
retirer.

\begin{notion}{Pourquoi le PID 1 n'est pas un processus comme un autre}
Le premier processus d'un conteneur hérite d'une obligation : adopter les
processus orphelins et récupérer leur code de sortie. Sans cela, les processus
terminés s'accumulent à l'état \emph{zombie} et finissent par saturer la table.

Un \cmd{bash} lancé comme PID 1 ne fait pas ce travail. \cmd{supervisord}, si.
\end{notion}

\role
L'entrypoint écrit une configuration adaptée au mode choisi, puis se
\textbf{remplace} par supervisord :

\begin{lstlisting}[language=bash]
exec "$(command -v supervisord)" -c "$CONF"
\end{lstlisting}

Le mot \cmd{exec} compte : sans lui, le shell resterait PID 1 et ne relaierait
pas les signaux. Un \cmd{docker stop} deviendrait un arrêt brutal au bout du
délai de grâce.

\apport
C'est le premier des trois niveaux de résilience du projet : si un service meurt
\emph{dans} la machine, il est relancé sans que l'utilisateur s'en aperçoive.

% -----------------------------------------------------------------------------
\outil{OpenSSH}{brique 18 · l'accès demandé par le sujet}

\quoi
Le serveur SSH : il donne un terminal distant chiffré, authentifié par mot de
passe ou par clé.

\pourquoi
Il n'y avait pas de choix à faire : c'est l'accès que le sujet impose, et c'est
le standard universel. La question intéressante est plutôt \emph{comment} on
l'expose.

Le conteneur écoute sur son propre port 22, et Docker le publie sur un port de
l'hôte dans la plage 8000--9000. C'est ce qui permet de faire tourner des
dizaines de serveurs SSH sur une seule machine : chacun croit avoir le port 22.

\role
Deux modes d'authentification. Par défaut, le mot de passe root généré. Mais si
l'utilisateur a enregistré sa clé publique sur son compte, la machine
\textbf{n'accepte plus que} cette clé --- l'authentification par mot de passe
est coupée.

\begin{piege}{Le port est publié avant que le service écoute}
Docker publie le port dès la création du conteneur, mais \cmd{sshd} met une
seconde ou deux à démarrer. Un test qui se connecte immédiatement échoue de
façon intermittente.

La correction n'est pas d'attendre « assez longtemps » --- une durée fixe est
toujours trop courte un jour --- mais de réessayer jusqu'à ce que le service
réponde, avec une limite.
\end{piege}

\apport
L'accès exigé par le cahier des charges, avec une option plus sûre que le mot
de passe pour qui le souhaite.

% -----------------------------------------------------------------------------
\outil{ttyd}{brique 19 · le terminal dans le navigateur}

\quoi
Un petit serveur qui expose un terminal dans une page web, par WebSocket.

\pourquoi
Le sujet ne le demande pas. Il répond à un problème réel : un utilisateur sur un
poste de salle de TP n'a pas toujours de client SSH, ni le droit d'en installer.

Les alternatives --- Shellinabox, abandonné ; Wetty, en Node.js donc une chaîne
d'outillage de plus --- ont été écartées. ttyd est un binaire unique, en C, de
quelques centaines de kilooctets.

\role
Lancé par supervisord sur le port 7681 du conteneur, publié sur l'hôte. Détail
qui compte : il ne lance pas un shell, il lance \cmd{login}.

\begin{lstlisting}[language=bash]
ttyd -p 7681 -W /bin/sh -c "cat /etc/issue; exec $(command -v login)"
\end{lstlisting}

Le port ne donne donc pas un accès ouvert : il présente une invite qui réclame
l'identifiant et le mot de passe.

\gros[0.8]{zoom-motd.png}{Ce que montre le port du terminal web : une invite de
connexion, pas un shell.}

\gros[0.8]{zoom-bienvenue.png}{Après authentification, le bandeau rappelle que
la machine est louée pour une durée limitée.}

\apport
Un accès sans rien installer. C'est le premier des ajouts au-delà du sujet.

% -----------------------------------------------------------------------------
\outil{Xvnc, XFCE et noVNC}{brique 20 · le bureau graphique}

\quoi
Trois pièces qui s'emboîtent. \textbf{Xvnc} est un serveur d'affichage qui
dessine dans une mémoire plutôt que sur un écran. \textbf{XFCE} est
l'environnement de bureau qui s'y affiche. \textbf{noVNC} est une passerelle qui
traduit le protocole VNC en WebSocket, pour qu'un navigateur puisse l'afficher.

\pourquoi
XFCE plutôt que GNOME ou KDE : il tourne correctement dans un gigaoctet de
mémoire, là où les autres en réclament deux ou trois. Pour une machine louée,
c'est la différence entre trois utilisateurs par serveur et dix.

noVNC plutôt qu'un client VNC à installer : même raisonnement que pour ttyd ---
l'utilisateur ne doit rien installer.

\role
Trois services de plus sous supervisord, en mode bureau seulement. Le serveur
d'affichage \textbf{n'écoute que sur la boucle locale} : on ne peut l'atteindre
qu'à travers la passerelle, qui est chiffrée quand un certificat est monté.

Le mot de passe VNC est limité à huit caractères par le protocole : ce sont les
huit premiers du mot de passe root, ce que l'interface explique.

\gros[0.86]{zoom-bureau.png}{Le bureau XFCE d'une machine Ubuntu louée, affiché
dans le navigateur. Cinq processus cohabitent dans ce conteneur.}

\apport
Le second ajout au-delà du sujet, et le plus spectaculaire en démonstration : un
bureau Linux complet, dans un onglet, sans rien installer.

% -----------------------------------------------------------------------------
\outil{Le script d'entrée}{brique 21 · ce qui assemble la machine}

\quoi
\fichier{1\_docker/entrypoint.sh}, 166 lignes de shell. C'est le programme que
Docker lance quand la machine démarre.

\pourquoi
Il fallait un endroit unique qui adapte la machine à ce que l'utilisateur a
demandé. Mettre cette logique dans le Dockerfile aurait obligé à six images
indépendantes ; la mettre dans l'application aurait dispersé le savoir.

\role
Trois tâches, dans l'ordre :

\begin{enumerate}
\item poser le mot de passe root et, s'il y en a une, la clé SSH --- auquel cas
  il \textbf{coupe} l'authentification par mot de passe ;
\item écrire la configuration supervisord correspondant au mode ;
\item se remplacer par supervisord.
\end{enumerate}

Il sert aussi d'outil : la plateforme l'appelle pour changer le mot de passe
d'une machine à chaud, sans la redémarrer.

\begin{lstlisting}[language=bash]
docker exec <conteneur> /usr/local/bin/entrypoint.sh setpass <nouveau>
\end{lstlisting}

\apport
Une seule recette pour six variantes, et la rotation du mot de passe sans
interruption de service.

% -----------------------------------------------------------------------------
\section*{Ce que la plateforme sait faire après la couche 3}

Elle fabrique de vraies machines Linux, accessibles de trois façons, durcies, et
qui se réparent toutes seules en interne. Mais quelqu'un doit encore décider
\emph{quand} les créer et \emph{quand} les détruire --- et sur quel serveur.
